%%
%% fall-lecture06.tex
%% 
%% Made by Alex Nelson <pqnelson@gmail.com>
%% Login   <alex@lisp>
%% 
%% Started on  2026-10-13T10:01:26-0700
%% Last update 2026-10-13T10:01:26-0700
%% 

\lecture{}

Our goal is to prove the following statement:

\begin{node}
Second-countable normal spaces are metrizable.
\end{node}

\begin{lemma}\label{lemma:fall:lec6:l1}
The Hilbert cube $H:=[0,1]^{\NN}$ is metrizable.
\end{lemma}

\begin{proof}[Proof sketch]
Elements of $H$ are sequences $\NN\to[0,1]$ in the interval $[0,1]$ so
$\vec{x}=(x(1),x(2),x(3),\dots)$ and each $x_{n}(j)\in[0,1]$. Then
$x_{n}$ converges pointwise to $x$ iff for every $j\in\NN$ we have
$x_{n}(j)\to x(j)$. Then we need to use the fact that $H$ is
first-countable: a neighborhood base at $x$ is given by
\begin{equation}
U_{n,m}:=\{y\mid |y(1)-x(1)|<\frac{1}{n},\dots,|y(m)-x(m)|<\frac{1}{n}\}
\end{equation}
are open sets and obviously $\{U_{m,n}\mid m,n\in\NN\}$ is countable.

Now we define a metric
\begin{equation}
\rho(x,y)=\max_{j}\frac{|x(j)-y(j)|}{2^{j}},
\end{equation}
and the proof (that this is a metric) is obvious.

Now we claim: $x_{n}\to x$ pointwise iff $\rho(x_{n},x)\to0$. Now, we
have two topologies on $H$ (the product topology and the metric
topology) which have the same set of convergent sequences, which means
they must be the same topology.
\end{proof}

\begin{lemma}\label{lemma:fall:lec6:l2}
Second-countable normal spaces have their closed sets be $G_{\delta}$.
\end{lemma}

\begin{proof}
Let $(X,\tau)$ be a second-countable normal space. Let $E\subset X$ be
a closed subset. Then $X\setminus E$ is second-countable. Therefore
$X\setminus E$ is Lindel\"{o}f in the relative topology by Theorem~\ref{thm:math110a:fall:second-countable-spaces-are-lindelof}.
For every $x\in X\setminus E$ we can separate $x$ and $E$ by disjoint
neighborhoods $x\in V$ and $E\subset U$ where $U\cap
V=\emptyset$. Then
\begin{subequations}
  \begin{align}
X\setminus E &= \bigcup_{x\in X\setminus E}V(x)\\
&=\bigcup_{n} V(x_{n})
  \end{align}
\end{subequations}
by Lindel\"{o}f there is a countable subcover (which is the second
line); then its complement
\begin{equation}
E = \bigcap_{n\in\NN}V(x_{n})^{\complement}\supset\bigcap_{n}U_{x_{n}}
\end{equation}
since $U_{x_{n}}\subset V(x_{n})^{\complement}$. But we also see that
\begin{equation}
E\subset\bigcap_{n}U_{x_{n}}
\end{equation}
since each $U_{x_{n}}$ is a neighborhood of $E$. Therefore we conclude
that the arbitrary closed subset is a countable intersection of open sets
\begin{equation}
E=\bigcap_{n\in\NN}U_{x_{n}}
\end{equation}
which is $G_{\delta}$.
\end{proof}

\begin{remark}
We can weaken the previous lemma from normal spaces $(T_{4})$ to
$T_{3}$ spaces.
\end{remark}

\begin{lemma}[Central point]\label{lemma:fall:lec6:l3}
If $(X,\tau)$ is normal and if $E\subset X$ is a closed $G_{\delta}$
set, then there exists a continuous function $f\colon X\to[0,1]$ such
that $f^{-1}(\{0\})=E$ is the zero set of $f$.
\end{lemma}

\begin{proof}
We see that 
\begin{equation}
X\setminus E=\bigcup_{n}F_{n}
\end{equation}
is the union of closed sets $F_{n}$ (because of the $G_{\delta}$ hypothesis),
and that for each $n$ we have:
\begin{equation}
F_{n}\cap E=\emptyset.
\end{equation}
Then by Urysohn's lemma, there are continuous functions $f_{n}\colon X\to[0,1]$
such that $f_{n}|_{E}=0$ and $f_{n}|_{F_{n}}=1$. Then we set
\begin{equation}
f(x) := \sum_{n}\frac{f_{n}(x)}{2^{n}},
\end{equation}
which is a continuous function, and $f|_{E}=0$ and $f|_{X\setminus E}>0$.
\end{proof}

\begin{theorem}[Urysohn metrizability]
Second-countable normal spaces are metrizable.
\end{theorem}

\begin{proof}
Let $(X,\tau)$ be a second-countable normal space. We will prove it is
metrizable. We have a countable base $\{U_{j}\}$. We write
$E_{j}:=X\setminus U_{j}$ which are closed $G_{\delta}$ sets by
Lemma~\ref{lemma:fall:lec6:l2}. Then by Lemma~\ref{lemma:fall:lec6:l3}
there are continuous functions $f_{j}\colon X\to[0,1]$ such that
$f^{-1}\{0\}=E_{j}$, and we combine the $f_{j}$'s to get the function
to the Hilbert cube:
\begin{equation}
f=(f_{1},f_{2},\dots)\colon X\to H.
\end{equation}
Then we claim:
\begin{enumerate}
\item $f$ is continuous (since for all $j$, $f_{j}=\pi_{j}\circ f$ and
  the definition of product topology),
\item $f$ is injective,
\item $f\colon X\to f(X)$ is a bijection.
\end{enumerate}
We need to prove the inverse map $f^{-1}$ is continuous to obtain a
homeomorphism onto its image (which is obvious). Then $f$ is a
homeomorphism onto its image, and its image $f(X)$ is metrizable,
therefore $X$ is metrizable.

\textsc{Claim 2: $f$ is injective}. Let $x\neq y$ be distinct
points. We will show $f(x)\neq f(y)$. We can separate $x$ and $y$:
there exists disjoint neighborhoods $V_{1}\ni x$ and $V_{2}\ni y$.
There exists base elements $U_{\ell}\subset V_{1}$ containing $x\in U_{\ell}$.
We write $E_{\ell}=X\setminus U_{\ell}$ which is a closed set,
therefore $G_{\delta}$. Then there exists an $f_{\ell}\colon X\to[0,1]$
where $f_{|ell}=0$ on $E_{\ell}$. This implies
$x\in U_{\ell}$, which implies $f_{\ell}(x)>0$, and together with
$y\notin U_{\ell}$ implies $f_{\ell}(y)=0$. Hence we obtain injectivity.

\textsc{Claim 3: $f$ is an open map}. Let $A\subset X$ be an open set.
We will show $f(A)$ is open in the relative topology in $f(X)$. It
suffices to show for any $j$ that $f(U_{j})$ is open in $f(X)$ (since
they cover $f(X)$). Then
\begin{equation}
f(U_{j})=f(X)\cap\{y\mid\pi_{j}y>0\}
\end{equation}
by the proof of claim 2, and it is the intersection of open sets with
$f(X)$, so it is relatively open. Hence the claim, hence the proof.
\end{proof}